% Shared body of the paper: everything the two drivers have in common.
% `main.tex` builds the named-author version, `main-anon.tex` the double-blind
% submission; they differ only in `\ifAnonymous`, which drives every
% `\Anon{...}{...}` below and in the section files.
%
% Do not add a \documentclass here -- the drivers own that.

\usepackage[utf8]{inputenc}

% !TEX root = main.tex
\usepackage[T1]{fontenc}
\usepackage{xcolor}
\usepackage{listings}
\usepackage{inconsolata}
\usepackage{relsize}
\usepackage{amsmath,amssymb,array}
\usepackage{esz}
\usepackage{mathpartir}
\usepackage{graphicx}
% The per-target tables are wider than the LNCS text block: set them in landscape.
\usepackage[figuresright]{rotating}
\usepackage{pgfplots}
\pgfplotsset{compat=1.18}
\usepgfplotslibrary{groupplots}
\usepackage{xspace}
\usepackage{xurl}
% Zapf Dingbats rather than bbding: bbding is a METAFONT bitmap font, which
% embeds as Type 3 with no Unicode map -- it extracts as garbage and arXiv
% flags it. \Checkmark/\XSolidBrush keep their names so call sites are
% unchanged; accsupp gives the glyphs readable text in the PDF's text layer.
\usepackage{pifont}
\usepackage{accsupp}
\DeclareRobustCommand{\Checkmark}{\BeginAccSupp{ActualText=ok}\ding{51}\EndAccSupp{}}
\DeclareRobustCommand{\XSolidBrush}{\BeginAccSupp{ActualText=failed}\ding{55}\EndAccSupp{}}
\usepackage{tabularx}
\usepackage{enumitem}
% Must load before `\lstMakeShortInline` makes `|` active, otherwise its own
% sources are read with that catcode and break.
\usepackage{booktabs}
\usepackage{microtype}

\usepackage{tcolorbox}
\tcbuselibrary{listings,breakable,skins}
% llncs fixes the page geometry; do not load geometry here.

% hypertexnames=false: llncs's chapter counter makes appendix sections reuse body anchors.
\usepackage[bookmarks,unicode,colorlinks=true,hypertexnames=false,
            linkcolor=blue,citecolor=blue,urlcolor=blue]{hyperref}
\def\UrlFont{\rmfamily}

% Inline code snippets are rigid boxes; stretch spaces rather than overrun the margin.
\setlength{\emergencystretch}{1.5em}

\newcommand{\Section}[1]{\section{#1}}
\newcommand{\SubSection}[1]{\subsection{#1}}
% LNCS third-level heading: bold, run-in, unnumbered; the class adds no period.
\newcommand{\Paragraph}[1]{\subsubsection{#1.}}

% Labeled-paragraph list: bold inline label, lightly indented.
% Use as \begin{Description} \item[Label.] body \item[...] ... \end{Description}.
\newlist{Description}{description}{1}
\setlist[Description]{leftmargin=1em, labelindent=1em, style=unboxed, font=\bfseries, itemsep=0.5ex, topsep=0.5ex}

% Bulleted list with an em-dash marker flush against the left text margin.
% Use as \begin{Itemize} \item ... \item ... \end{Itemize}.
\newlist{Itemize}{itemize}{1}
\setlist[Itemize]{label=--, leftmargin=1em, labelindent=0pt, labelsep=0.5em, itemsep=0.5ex, topsep=0.5ex}

% Double-blind support: the drivers declare `\ifAnonymous`, set only by
% main-anon.tex, so the same sources serve both.
% \Anon{normal}{anonymized} -- use it for anything naming the authors,
% their employer, or their internal systems.
\newcommand{\Anon}[2]{\ifAnonymous #2\else #1\fi}

\newcommand{\TODO}[2]{%
    \noindent\textcolor{gray}{%
            {\scriptsize\textbf{TODO(#1)}: #2}
    }
}

\newenvironment{Figure}{
    \begin{figure}[tb]
} {
    \end{figure}
}

% Full-width variant; in the one-column LNCS layout figure* is identical to
% figure.
\newenvironment{WideFigure}{
    \begin{figure*}[!t]
} {
    \end{figure*}
}

% Full-width table variant, see WideFigure.
\newenvironment{WideTable}{
    \begin{table*}[!t]
} {
    \end{table*}
}

% Abbreviations

\newcommand{\MVP}[0]{\textsf{MVP}\xspace}
\newcommand{\MSL}[0]{\textsf{MSL}\xspace}
\newcommand{\WP}[0]{WP\xspace}
\newcommand{\solidity}[0]{\textsf{Solidity}\xspace}
\newcommand{\kay}[0]{\textsf{K}\xspace}
\newcommand{\fstar}[0]{\textsf{F*}\xspace}
\newcommand{\verus}[0]{\textsf{Verus}\xspace}
\newcommand{\certora}[0]{\textsf{Certora}\xspace}
\newcommand{\CC}[0]{\textsf{CC}\xspace}


% IVL

% Based on esz.sty -- see CLAUDE.md

\newenvironment{ivl}{\zedindent=2ex\begin{zed}}{\end{zed}}

\newcommand{\requiresof}[2]{\Zpreop{requires\_of}\langle#1\rangle(#2)}
\newcommand{\abortsof}[2]{\Zpreop{aborts\_of}\langle#1\rangle(#2)}
\newcommand{\ensuresof}[2]{\Zpreop{ensures\_of}\langle#1\rangle(#2)}
\newcommand{\resultof}[2]{\Zpreop{result\_of}\langle#1\rangle(#2)}
\newcommand{\modifiesof}[1]{\Zpreop{modifies\_of}\langle#1\rangle}

% Variant-set names used in the implementation section.
\newcommand{\Clo}{\mathcal{C}}
\newcommand{\Par}{\mathcal{P}}
\newcommand{\Fld}{\mathcal{F}}

\newcommand{\Addr}{\Zpreop{address}}

\newcommand{\Mem}{\mathcal{M}}
\newcommand{\Mod}{\Delta}
\newcommand{\Type}{\mathcal{T}}
\newcommand{\Expr}{\mathcal{E}}

% Keywords for  IVL pseudocode
% ... esz already has \LET \IF \THEN \ELSE
\newcommand{\PROC}{\Zkeyword{proc}}
\newcommand{\FUN}{\Zkeyword{fun}}
\newcommand{\AXIOM}{\Zkeyword{axiom}}
\newcommand{\TRIGGER}{\Zinrel{trigger}}
\newcommand{\RETURNS}{\Zinrel{returns}}
\newcommand{\HAVOC}{\Zkeyword{havoc}}
\newcommand{\ASSUME}{\Zkeyword{assume}}
\newcommand{\ASSERT}{\Zkeyword{assert}}
\newcommand{\DATATYPE}{\Zkeyword{datatype}}
\newcommand{\IS}{\Zinrel{is}}
\newcommand{\TRUE}{\Zkeyword{true}}
\newcommand{\OLD}{\Zpreop{old}}
\newcommand{\BOOL}{\Zpreop{bool}}
\newcommand{\READ}{\Zpreop{read}}
\newcommand{\CALL}{\Zpreop{call}}
\newcommand{\VAR}{\Zkeyword{var}}
\newcommand{\FOREACH}{\Zkeyword{foreach}}

% State-labeled expression: \slabeled{S}{e} renders as S \sat e.
% \sat is a bar-then-tilde composite that reads like the literal |~.
% No standard single-glyph relation in amssymb / stmaryrd matches
% (closest is \vDash, which is |=). We therefore compose: \scriptstyle
% shrinks both glyphs; -1mu tightens the tilde against the bar while
% keeping a small visible gap.
% Kept in sync with the Move listings literate mapping below.
\newcommand{\sat}{\mathrel{{\scriptstyle\vert}\mkern-1mu{\scriptstyle\sim}}}
\newcommand{\slabeled}[2]{#1 \sat #2}

\newcommand{\sq}[1]{\overline{#1}}


% Source Transformation

\newcommand{\transform}[0]{\Large{$\leadsto$}}


%%%% Code

% Keep the LNCS default font; do not override with charter.
%\usepackage[charter]{mathdesign}
%\def\rmdefault{bch}
%\def\ttdefault{blg}

\lstdefinestyle{MoveStyle}{
    basicstyle=\ttfamily,
    keywordstyle=\color{black}, % don't bold keywords which is the default
    commentstyle=\color{gray}\normalfont\itshape,
    escapechar=@, % use to embed LaTeX into code
    breaklines=true,
    % \sat is set in \scriptstyle and followed by tight kerning, so its
    % right side-bearing is small; tail the replacement with \, so the
    % symbol doesn't abut the next token visually.
    literate={|~}{{$\sat$\,}}1
        {<=}{{$\leq$}}2
        {>=}{{$\geq$}}2
        {==}{{$=$}}2
        {!=}{{$\neq$}}2
        {==>}{{$\Longrightarrow$}}3
        {forall}{{$\forall$}}1
        {exists}{{$\exists$}}1
        % Later entries win: `exists<R>(a)` is the resource predicate, not a quantifier.
        {exists<}{{exists<}}7,
}


\lstdefinelanguage{Move}{
    morekeywords={
        abort,
        aborts_if,
        aborts_of,
        acquires,
        address,
        apply,
        as,
        assert,
        assume,
        borrow_global,
        borrow_global_mut,
        break,
        const,
        continue,
        copy,
        copyable,
        define,
        drop,
        else,
        ensures,
        ensures_of,
        exists,
        false,
        forall,
        friend,
        fun,
        global,
        has,
        havoc,
        if,
        in,
        include,
        invariant,
        key,
        let,
        loop,
        match,
        modifies,
        modifies_of,
        module,
        move,
        move_from,
        move_to,
        mut,
        native,
        num,
        old,
        onabort,
        pragma,
        proof,
        public,
        requires,
        requires_of,
        resource,
        result_of,
        return,
        schema,
        script,
        signer,
        spec,
        split,
        store,
        struct,
        true,
        u8,
        u64,
        u128,
        u256,
        update,
        use,
        with,
        where,
        while},
    sensitive=true,
    morecomment=[l]{//},
    morecomment=[s]{/*}{*/},
}

% Short inline delimiters for Move code: both |...| and !...! work.
% The ! alternative is useful when the snippet itself contains |,
% e.g. function-type syntax !|T1,...,Tn| -> T!.
%
% Inline snippets must never wrap mid-token (e.g. between & and T in |&T|),
% but they may still break at whitespace so a long inline snippet does not
% overflow the column. Block listings (Move environment) keep MoveStyle's
% default break-anywhere behavior.
\lstMakeShortInline[language=Move,style=MoveStyle,breaklines=true,breakatwhitespace=true,escapechar=@]|
\lstMakeShortInline[language=Move,style=MoveStyle,breaklines=true,breakatwhitespace=true,escapechar=@]!


% This defines a new environment for Move code.
\lstnewenvironment{Move}{
    \lstset{
        language=Move,
        style=MoveStyle,
        basicstyle=\relsize{-1}\ttfamily,
        keywordstyle=\color{blue},
    }
}{
}

% This defines a new environment for Move code in a box, without line numbers.
\lstnewenvironment{MoveBox}{
    \lstset{
        language=Move,
        style=MoveStyle,
        basicstyle=\relsize{-1}\ttfamily,
        keywordstyle=\color{blue},
    %numbers=left,
    %numberstyle=\scriptsize\color{gray},
    %frame=single,
    }
}{
}

% Provenance-tinted variants of MoveBox: `MoveBoxUser` for the program the
% user brings, `MoveBoxAgent` for text the coding agent wrote, `MoveBoxTool`
% for verbatim output of an MCP tool, which it names (`{wp}`, `{verify}`).
%
% tcolorbox rather than a listings frame: the tag is overlaid on the first line
% of the box, and listings paints its own `backgroundcolor` per line *after*
% anything emitted before it, which would bury the tag.
% Per-cell dot plot: one hue per model, chrome in neutral grays.
\definecolor{ModelSol}{HTML}{2A78D6}
\definecolor{ModelTerra}{HTML}{EB6834}
\definecolor{CellMuted}{HTML}{898781}
\definecolor{CellGrid}{HTML}{E1E0D9}
\definecolor{CellAxis}{HTML}{C3C2B7}

\definecolor{UserRule}{HTML}{3F8A4E}
\definecolor{UserBack}{HTML}{F2F9F3}
\definecolor{AgentRule}{HTML}{2C6FB5}
\definecolor{AgentBack}{HTML}{F2F7FC}
\definecolor{ToolRule}{HTML}{B07A1E}
\definecolor{ToolBack}{HTML}{FCF8EF}

\tcbset{
    ProvenanceBox/.style={
        % `enhanced` is required for `overlay`: the plain `standard` skin
        % ignores it, so an unbroken box would silently lose its tag.
        enhanced,
        listing only,
        listing options={
            language=Move,
            style=MoveStyle,
            basicstyle=\relsize{-1}\ttfamily,
            keywordstyle=\color{blue},
        },
        breakable,
        sharp corners,
        boxrule=0pt,
        leftrule=1.5pt,
        boxsep=0pt,
        left=6pt, right=3pt, top=3pt, bottom=3pt,
    },
    % Tag only on the first piece of a box that breaks across pages.
    ProvenanceTag/.style 2 args={
        overlay unbroken and first={%
            \node[anchor=north east, inner xsep=4pt, inner ysep=2pt,
                  font=\scriptsize\sffamily, text=#1] at (frame.north east) {#2};},
    },
}

\newtcblisting{MoveBoxUser}{
    ProvenanceBox,
    colback=UserBack, colframe=UserRule,
    ProvenanceTag={UserRule}{user},
}

\newtcblisting{MoveBoxAgent}{
    ProvenanceBox,
    colback=AgentBack, colframe=AgentRule,
    ProvenanceTag={AgentRule}{agent},
}

\newtcblisting{MoveBoxTool}[1]{
    ProvenanceBox,
    colback=ToolBack, colframe=ToolRule,
    ProvenanceTag={ToolRule}{#1 tool},
}

% This defines a new environment for Move code in a box, with line numbers.
\lstnewenvironment{MoveBoxNumbered}{
    \lstset{
        language=Move,
        style=MoveStyle,
        basicstyle=\relsize{-1}\ttfamily,
        keywordstyle=\color{blue},
        numbers=left,
        numberstyle=\scriptsize\color{gray},
    %frame=single,
    }
}{
}

% This defines a new environment for diagnostics as produced by the prover.
\lstnewenvironment{MoveDiag}{
    \lstset{
        style=MoveStyle,
        basicstyle=\relsize{-2}\ttfamily,
    }
}{
}

% Environment for Claude Code transcripts.
%
% Source convention: lines that begin with "> " are the user prompt;
% everything else is the model's response. The "> " is rendered as "❯ "
% in the PDF (via a literate substitution) so the output matches the
% Claude Code CLI's prompt glyph without forcing the author to type a
% non-ASCII character.
%
% Prompt lines are styled bold and blue so they stand out; response lines
% are rendered in plain monospace, mimicking how the transcript looks in
% a terminal.
%
% Unicode glyphs commonly emitted by Claude Code (tree connectors,
% checkboxes, ellipsis, arrows, check/cross marks) are mapped via the
% literate= option so they render correctly under pdfLaTeX (where listings
% does not handle multi-byte input directly).
%
% Wrapping: long prompts wrap with breakindent matched to the width of
% "❯ " in monospace, so continuation lines line up under the prompt text
% (approximately --- listings uses a fixed indent, not a measured prefix).
%
% Limitation: ">" at the start of a line is reserved as the prompt marker.
% A literal ">" inside a response line must be escaped via the listings
% escapechar (@), e.g. @\textgreater @.
\lstnewenvironment{claude}{
    \lstset{
        basicstyle=\relsize{-1}\ttfamily,
        breaklines=true,
        breakautoindent=true,
        breakatwhitespace=true,
        breakindent=1.2em,
        columns=fullflexible,
        keepspaces=true,
        escapechar=@,
        extendedchars=true,
        inputencoding=utf8,
        literate=
            {⎿}{{$\llcorner$}}1
            {└}{{$\llcorner$}}1
            {├}{{$\vdash$}}1
            {─}{{-}}1
            {│}{{$\vert$}}1
            {◻}{{$\square$}}1
            {◼}{{$\blacksquare$}}1
            {☐}{{$\square$}}1
            {☑}{{$\boxtimes$}}1
            {●}{{$\bullet$}}1
            {○}{{$\circ$}}1
            {◯}{{$\bigcirc$}}1
            {⏺}{{$\bullet$}}1
            {✽}{{$\ast$}}1
            {★}{{$\star$}}1
            {☆}{{$\star$}}1
            {…}{{\ldots}}1
            {↗}{{$\nearrow$}}1
            {↘}{{$\searrow$}}1
            {↳}{{$\hookrightarrow$}}1
            {✓}{{\Checkmark}}1
            {✗}{{\XSolidBrush}}1
            {•}{{$\bullet$}}1
            {–}{{--}}1
            {—}{{---}}1
            {‘}{{`}}1
            {’}{{'}}1
            {“}{{``}}1
            {”}{{''}}1,
        moredelim=[l][\color{blue}\bfseries]{>},
    }
}{
}


%%% Local Variables:
%%% mode: latex
%%% TeX-master: "main"
%%% End:

% Generated by eval_cells.py; do not edit.
\newcommand{\CellTaskLabels}{BA-base-012,BK-bucket-016,LP-price-021,MD-median-015,MM-min-013,OV-order-006,PM-curve-027,QP-part-025,SM-select-022,TL-lev-020,TR-cancel-026,TR-discard-011,TR-order-010,TS-trial-019,UC-credits-008,VS-contrib-003,VS-fees-001,VS-redeem-004,VS-shares-002,WU-consume-023}


\bibliographystyle{splncs04}

\title{Hybrid Mechanical and Agentic\texorpdfstring{\\}{} Specification Inference}
\titlerunning{Hybrid Mechanical and Agentic Specification Inference}

\Anon{%
	\author{Wolfgang Grieskamp \and Teng Zhang \and Vineeth Kashyap}
	\authorrunning{W. Grieskamp et al.}
	\institute{Aptos Labs, Palo Alto, USA\\ \email{wg@aptoslabs.com}}
}{%
	\author{Anonymous Author(s)}
	\institute{Anonymous Institution(s)}
}

\begin{document}

\maketitle

\begin{abstract}
Formal verification suffers from the cost of specification: to establish high-level facts, contracts have to be written for supporting functions, many of them mechanical and low-level. 
Specification inference derives such contracts from the program. 
We present a hybrid agentic-mechanical inference approach for the Move smart contract language, which has many similaritiers with Rust, and is integrated with the Move Prover.
An agent controls inference and verification via a tool which offers, beside calling the Move Prover, deriving specs from code via weakest precondition analysis. 
Equipped with this tool, and loop invariants guessed by AI, the agent is capable to derive complete specs for non-trivial problems.
We evaluate the approach on 20 verification targets extracted from a private production codebase, with completeness checked by withheld mutants, under a frontier model, Sol~5.6, and a mid-tier model of the same family, Terra~5.6. The unaided agent completes 93\% of its cells. With \WP, both models complete more, and the mid-tier model completes as many as the frontier model at 55\% of its cost.

	\keywords{Specification inference \and Weakest precondition \and Reference elimination \and LLM agents \and Move Prover}
\end{abstract}

% !TEX root = main.tex
\Section{Introduction}
\label{sec:intro}

A roadblock in program verification adoption is the \emph{cost of specification}.
The properties developers care about are high-level --- a call chain aborts only under stated preconditions, an entrypoint satisfies its postconditions, stored data respects an invariant --- but to prove any of them, the supporting functions need contracts of their own, and those are often boilerplate, as they restate what the code does.
\emph{Specification inference} derives such contracts from the program, so that effort can go to the properties worth stating by hand.
The classical engine for mechanical inference is Dijkstra's weakest-precondition transformation (\WP)~\cite{dijkstra1975wp}: propagated backward through a function, it yields the exact abort condition and the exact postcondition. However, it depends on loop invariants to provide an induction hypothesis~
\cite{BL05}.  Invariant synthesis has no complete technique despite a long history~\cite{CousotHalbwachs78} of research.

This paper describes a novel combination of agentic specification inference -- researched e.g.~in \cite{INF-PAPERS} -- and weakest precondition transformation.
The agent conducts the inference tactic, supplies loop invariants, and calls \WP and verification as a tool to aide the work.
The approach is implemented for the Move smart contract language which shares many simularities with Rust. Move is equipped with strong verification support by the Move Prover (\MVP)~\cite{DillGPQXZ22,GrieskampEtAl26HO}.
Move shares the reference logic with safe Rust, and is therfore amenable to reference elimination, which has been applied in the Move Prover since its incubation~\cite{DillGPQXZ22}. 
On reference-eliminated Move IR, we run classical \WP, mapping Moves reference logic into the predicate calculus of the spec language. 
Moreover, via the concept of behavioral predicates \cite{FMCAD}, which allow to reason about the specifcation of others functions as a black box, \WP becomes fully modular.

The combination is evaluated on a corpus of 20 actual inference problems drawn from a private codebase (\Anon{that of the perpetual trading system Decibel~\cite{Decibel}}{that of a production perpetual trading system}) --- private so that contracts must be inferred rather than recalled --- with three tactics, one without the \WP tool and two with, under two models of one family, the frontier model Sol~5.6 and the mid-tier model Terra~5.6, and with completeness checked by withheld mutants.
Without the \WP tool the agent already does well: it produced a complete specification in 93\% of its cells under both models, including on targets that require inventing lemmas and recursive helper functions.
With \WP, both models complete more cells, and the remaining failures occur on one target whose contract needs a permutation argument that \WP cannot supply either.
The effect on cost depends on the model.
For the frontier model, the hybrid tactics cost about the same as the agent alone, since the model finds the invariants unaided and then carries the derived contract in its context.
For the mid-tier model, they cost 20\% less, and with the tool it completes as many cells as the frontier model at 55\% of its cost.

\Paragraph{Contributions}
The paper presents a novel combination of agentic and mechanic specification inference, integrated with the Move Prover, with a detailed empirical evaluation on a set of real-world code examples, demonstrating that combinations like this can lower the cost and increase realiability in this domain. As a side-effect, we demonstrate that reference elimination in the style of the Move Prover is a good starting point for specifcation inference.

The Move Prover was developed for the Diem project and is \Anon{maintained at Aptos Labs since 2022}{since maintained by the blockchain company behind it}; it has verified Diem's payment system and critical components of \Anon{the Aptos blockchain}{a public blockchain}~\cite{ParkZGXGCLC24,AptosFrameworkBook}.
Move shares Rust's safe references, linear types, algebraic data types, and higher-order functions~\cite{MoveOnAptosBook}; the mechanical half of our approach relies on exactly these traits, so we expect it to carry over to verifiers for Rust-like languages that eliminate references the same way.

\endinput

% !TEX root = main.tex
\Section{Hybrid Inference by Example}
\label{sec:example}

Hybrid inference combines two sources of specifications: the mechanical \WP analysis of Sec.~\ref{sec:wp}, which derives exact contracts but needs loop invariants to do so, and a coding agent, which supplies the loop invariants and whatever else the analysis cannot derive, and repairs the specification based on verifier feedback.
The prover is connected to the agent through an MCP service that offers a tool for verification and a tool for \WP inference, and the agent carries a set of skills describing the language, the tools, and the workflow (Sec.~\ref{sec:skills}).
We walk through one session; the example is written in Move but should be readable without knowing the language~\cite{MoveOnAptosBook}.

\label{sec:pow}
The example computes the power $\mathit{base}^{\mathit{exp}}$ iteratively:

\begin{MoveBoxUser}
  fun pow(base: u64, exp: u64): u64 {
    let result = 1;
    let i = 0;
    while (i < exp) {
      result = result * base;
      i = i + 1;
    }
    result
  }
\end{MoveBoxUser}

Following its skill instructions, the agent first calls \WP. The answer reports the missing loop invariant and, as evidence, unrolls the loop symbolically and shows the values of the loop variables at each of the first loop entries:

\begin{MoveBoxTool}{wp}
WP inferred `vacuous` conditions after this loop without an invariant. ..
  --> sources/pow.move:5:14
  |
5 |       while (i < exp) {
  |              ^
  |
  = loop-invariant evidence 
  = bounded loop-head facts :
      head[0]: head[0].i == 0
               head[0].result == 1
      head[1]: 0 < exp ==> head[1].i == 1
               0 < exp ==> head[1].result == base
      head[2]: 1 < exp ==> head[2].i == 2
               1 < exp ==> head[2].result == base * base
      ..
\end{MoveBoxTool}

From this evidence, the agent derives the loop invariant and introduces a recursive specification helper function for the accumulated result.
It then calls the verification tool to check whether the invariant verifies:

\begin{MoveBoxAgent}
  fun pow(base: u64, exp: u64): u64 {
    ..
    while (i < exp) {
      result = result * base;
      i = i + 1;
    } spec {
      invariant i <= exp;
      invariant result == math_pow(base, i);
    }
    ..   
  }
  // specification helper function
  spec fun math_pow(base: num, exp: num): num {
    if (exp <= 0) 1
    else math_pow(base, exp - 1) * base
  }
\end{MoveBoxAgent}

In the next turn the agent runs the \WP tool again.
With the loop invariant present, \WP derives the complete specification.
However, it marks the abort condition as |sathard|. The marker is set when a derived condition is exact but its satisfiability is expected to be hard for the solver --- here because of the existential quantifier, which states that the computation aborts if the multiplication overflows in some iteration |y|.
The specification language of the Move Prover uses unbounded integers, |num|, which makes it possible to state overflow of the program's |u64| values:

\begin{MoveBoxTool}{wp}
  spec pow { // post and aborts condition for pow
    ensures result == math_pow(base, exp);
    aborts_if [inferred = sathard] exists y: u64:
      math_pow(base, y) * base > MAX_U64
      && (forall x: num: 0 <= x && x <= y
            ==> math_pow(base, x) <= MAX_U64)
      && y < exp
  }
\end{MoveBoxTool}

The agent tries to verify this specification, but verification times out with the following diagnosis:

\begin{MoveBoxTool}{verify}
  verification failed:
  error: verification out of resources/timeout
         (timeout set to 5s)
   = timeout analysis replay:
       mixed quantifier and nonlinear arithmetic
   = top quantifier activity:
       quantifier at sources/pow.move:24
         -- 55,793+ observed instantiations
\end{MoveBoxTool}

The quantifier the diagnosis points to is the recursive definition of the helper function |math_pow|, which the solver keeps unfolding while it tries to prove the |sathard| abort condition. In the next turn, the agent introduces lemmas that bound every partial product by the last one, which lets it restate the abort condition without a quantifier:

\begin{MoveBoxAgent}
  spec pow {
    pragma opaque;
    ensures result == math_pow(base, exp);
    /// The partial products are non-decreasing, so the
    /// loop overflows exactly when its last product does.
    aborts_if
      exp > 0 && math_pow(base, exp - 1) * base > MAX_U64;
  } proof {
    // Bounds every iteration's product by the last one;
    // triggers on the invariant's `math_pow(base, i)`.
    forall k: num {math_pow(base, k)}
      apply math_pow_last_is_max(base, k, exp);
  }
  spec lemma math_pow_last_is_max(
      base: num, k: num, n: num) {
    requires base >= 0 && 0 <= k && k < n;
    ensures math_pow(base, k) * base
              <= math_pow(base, n - 1) * base;
  } proof {
    apply math_pow_mono(base, k + 1, n);
  }
  spec lemma math_pow_mono ...  // left out for brevity
\end{MoveBoxAgent}

\noindent This specification passes verification, and the agent returns it as the result.

\endinput

% !TEX root = main.tex
\Section{Weakest Preconditions over Eliminated References}
\label{sec:wp}

The mechanical half of the system is a weakest-precondition analysis~\cite{dijkstra1975wp} that reads off abort conditions, modifies clauses, and consequence-style postconditions from a function's bytecode.
\WP itself is textbook.
Two things make it applicable to a language with Rust-style references and make its output readable: it runs on \MVP's reference-eliminated bytecode, where backward predicate propagation is plain substitution, and a reverse mapping carries the derived predicates back to the references of the source.
This section describes both, and then what the elimination buys compared to reasoning over references directly.

\SubSection{Reference Elimination}
\label{sec:refelim}

Move has Rust-style references: |&T| for shared reads and |&mut T| for exclusive writes, both with a stack-bounded lifetime, and the borrow checker guarantees that a live |&mut| is the only access path to its referent.
\MVP exploits this guarantee to rewrite bytecode into a reference-free form~\cite{DillGPQXZ22}, originally in order to keep the verification conditions first-order and free of a heap model.
Two passes perform the elimination.
Shared references are inlined: |&T| disappears and reads through it become reads of the underlying value.
A mutable reference becomes a pair of the value it currently denotes and a \emph{borrow path} back to its root --- a local, a |&mut| parameter, or a global resource |R[a]| --- and every mutation is followed by an explicit \emph{write-back} along that path, inserted by a borrow analysis, which rebuilds the parent value with the field or element replaced.
A |&mut| parameter becomes an in/out value parameter: it enters as a plain value and leaves as the value written back last.
Where the root of a reference is not statically unique, as in |if (c) &mut x else &mut y|, the write-back is guarded by an explicit test of which parent the reference was taken from.
Before \WP runs, the bytecode further passes through \MVP's specification instrumentation: each function-level |requires|, |aborts_if|, |ensures|, or |modifies| clause and each loop |invariant| becomes an |assume| or |assert| at the appropriate program point, and calls are summarized by the \emph{behavioral predicates} of~\cite{GrieskampEtAl26HO}: |requires_of<f>|, |aborts_of<f>|, |ensures_of<f>|, |modifies_of<f>|, and |result_of<f>| reify a callee's contract as predicates over the call site's arguments and result.
At the input to \WP, the bytecode is a sequence of assignments over plain values interleaved with |assume| and |assert|.

\SubSection{The Predicate Transformer}
\label{sec:transformer}

The analysis is a backward dataflow pass over the control-flow graph.
Its state at a program point is a triple $(E, A, \ell)$: a set $E$ of \emph{ensures} conjuncts that hold on every normal return reached from the point, a set $A$ of \emph{abort} disjuncts under which the function aborts, and a \emph{memory label} $\ell$ naming the memory version the conditions refer to, together with bookkeeping for the |&mut| parameters and global resources that have been written on the way.
At a return, $E = \{\,\mathtt{result}_i = t_i\,\}$ for the returned temporaries and $A = \emptyset$; at an |abort|, $A = \{\mathtt{true}\}$.
Instructions transform the state as follows.
\begin{Itemize}
\item \emph{Assignment and operators.}
  $t := e$ substitutes $e$ for $t$ in $E$ and $A$; an operator that can abort --- arithmetic overflow and underflow, division by zero, an index out of range, a cast that truncates, a missing enum variant --- also adds its abort predicate to $A$.
\item \emph{References.}
  A borrow $r := \,$|&x| or $r := \,$|&x.f| substitutes |x| or |x.f| for $r$; a write-back along a field edge replaces the parent |x| by |x| with field |f| set to the value written; a global borrow $r := \,$|&mut R[a]| substitutes |R[a]| for $r$ and adds |!exists<R>(a)| to $A$.
  Since every reference is a path to a named root, no instruction needs a heap.
\item \emph{Calls.}
  A call $t := f(\bar a)$ substitutes |result_of<f>|$(\bar a)$ for $t$ --- or the callee's specification function, if $f$ is pure --- adds |aborts_of<f>|$(\bar a)$ to $A$, and relates the memory before and after the call through |ensures_of<f>|; when the callee has no abort contract, the abort condition is recorded as \emph{partial} rather than invented.
  Calls are thus summarized modularly: the caller's contract says what $f$ returns without committing to how.
\item \emph{Branches.}
  At a branch on $c$ the two incoming states are merged path-sensitively: conditions common to both sides are kept as they are, the others are guarded, $c \Rightarrow e$ and $\neg c \Rightarrow e'$ in $E$, $c \wedge a$ and $\neg c \wedge a'$ in $A$.
\item \emph{Specifications.}
  An |assert| $\varphi$ adds $\varphi$ to $E$ and $\neg\varphi$ to $A$; an |assume| $\varphi$ adds $\varphi$ as an antecedent to both.
  This is how the contracts of callees, user-written conditions, and loop invariants enter the derivation.
\item \emph{Loops.}
  Loops are broken at their headers by the cut-point encoding of Barnett and Leino~\cite{BL05}: the back-edge is severed, the loop |invariant| is asserted before and after each iteration, and the loop-modified locations are havoced in between.
  \WP of |havoc x| universally quantifies |x| in $E$ and existentially in $A$, so variables that the invariant does not constrain appear quantified in the contract.
\end{Itemize}
The pass iterates to a fixpoint in reverse topological order of the acyclic graph that remains, and the state at entry is the raw contract.
Post-processing then strips the labels of the entry and exit memories, resolves the parent tests of non-unique write-backs into path conditions computed from the dominator tree, adds |p == old(p)| for every |&mut| parameter never written, substitutes the global resources captured through borrows, and simplifies: constant folding, boolean and arithmetic identities, tautology removal, and scoping of quantifiers to the conjuncts that mention their variable.
Simplification matters more here than in a verifier, where the \WP result goes to the solver and is never shown: the raw predicate of a twenty-line function can have hundreds of terms, and the agent and the user see the simplified one.

\SubSection{Carrying the Result Back}
\label{sec:reverse}

The predicate at entry is stated over the eliminated form, whose in/out parameters and rebuilt parent values have no counterpart in the source signature; the contract the user reads must be stated over the original references.
Three cases arise, and each has a source-level spelling.
A |&T| parameter needs nothing: the specification language already reads a shared reference as its value.
For a |&mut T| parameter |p|, the in/out parameter's exit value becomes |ensures p == e|, where |e| is the value written back last on the path to the return, and any read of |p| before that write resolves to |old(p)|, the value at entry; an earlier write is folded in by substituting it for |old(p)| in what was already derived, so a sequence of writes chains into one equation.
For a global resource, a borrow of |R[a]| that is written back yields |modifies R[a]| together with an equation on the resource's post-state value, with |old(R[a])| for its value at entry.
Consider the transfer of an amount between two balances:

\begin{MoveBoxUser}
  struct Balance has key { v: u64 }
  fun transfer(from: address, to: address, amount: u64) {
    Balance[from].v -= amount;
    Balance[to].v += amount;
  }
\end{MoveBoxUser}

\noindent Both statements borrow a resource mutably, read and write a field through the reference, and write back.
The derived contract is stated over the resources themselves, and the two borrow roots may coincide, so \WP keeps them apart by a case split on |from == to|:

\begin{MoveBoxTool}{wp}
  spec transfer {
    pragma opaque;
    modifies Balance[from];
    modifies Balance[to];
    aborts_if [inferred] !exists<Balance>(from);
    aborts_if [inferred] !exists<Balance>(to);
    aborts_if [inferred] Balance[from].v < amount;
    aborts_if [inferred]
      from != to && Balance[to].v + amount > MAX_U64;
    ensures [inferred]
      from == to ==> Balance[from] == old(Balance[from]);
    ensures [inferred] from != to
      ==> Balance[from].v == old(Balance[from]).v - amount;
    ensures [inferred] from != to
      ==> Balance[to].v == old(Balance[to]).v + amount;
  }
\end{MoveBoxTool}

\noindent When a function updates global memory in several steps, an equation on the final memory is not enough: the abort condition of the second |move_to| below depends on the memory after the first.
The memory labels of the transformer therefore surface as \emph{state labels} naming the intermediate memory states at which a global update happens, and the specification language lets a condition be evaluated at a labeled state, !S |~ e!, or over a range of states, !S1..S2 |~ e!:

\begin{MoveBoxUser}
  fun split_balance(
      dst1: &signer, dst2: &signer, src: address): u64 {
    let Balance { v } = move_from<Balance>(src);
    let half = v / 2;
    move_to(dst1, Balance { v: half });
    move_to(dst2, Balance { v: v - half });
    half
  }
\end{MoveBoxUser}
\begin{MoveBoxTool}{wp}
  spec split_balance {
    pragma opaque;
    modifies Balance[src];
    modifies Balance[signer::address_of(dst1)];
    modifies Balance[signer::address_of(dst2)];
    ensures [inferred] result == old(Balance[src]).v / 2;
    ensures [inferred] ..S1 |~ remove<Balance>(src);
    ensures [inferred] S1..S2 |~
      publish<Balance>(signer::address_of(dst1),
        Balance { v: old(Balance[src]).v / 2 });
    ensures [inferred] S2.. |~
      publish<Balance>(signer::address_of(dst2),
        Balance { v: old(Balance[src]).v
                     - old(Balance[src]).v / 2 });
    aborts_if [inferred] !exists<Balance>(src);
    aborts_if [inferred]
      S1 |~ exists<Balance>(signer::address_of(dst1));
    aborts_if [inferred]
      S2 |~ exists<Balance>(signer::address_of(dst2));
  }
\end{MoveBoxTool}

\noindent |S1| names the memory between the |move_from| and the first |move_to|, |S2| the memory between the two |move_to|s; |publish<R>(a, x)| states that resource |R| at |a| equals |x| over the range, |remove<R>(a)| that none is present, and |exists<R>(a)| is the usual existence predicate.
Without the labels, the memory versions the eliminated form threads through the function could be read back only as a single, weaker equation on the final memory.

\SubSection{Diagnostics for the Agent}
\label{sec:diagnostics}

\WP is complete only where its inputs are: a loop without an invariant and a callee without an abort contract both leave holes, and the tool reports both instead of silently emitting an incomplete contract.
After a havoced loop, a condition that quantifies over everything the loop touched carries no information; such \emph{vacuous} conditions are dropped and the loop is reported with \emph{loop-invariant evidence} --- the facts that hold at the first few loop heads on each path, computed by bounded unrolling of the same transformer, as in the |head[k]| listing of Sec.~\ref{sec:example}.
A condition that is exact but existentially quantified, typically an abort condition that says \emph{some} iteration overflows, is kept and marked |sathard|, since its satisfiability will be expensive for the solver.
An abort condition that could not account for a callee is marked \emph{partial}.
Every emitted condition carries an |[inferred]| marker, so that re-runs replace their own output and never disturb user-written specifications, and every function that received a contract is made |opaque|, so that callers reason through it.
The agent's skills are written against these markers: they identify the loop that needs an invariant, the condition to simplify, and the callee to specify first.

\SubSection{What the Elimination Buys}
\label{sec:advantages}

Reference elimination was introduced to make \emph{verification} of Move tractable~\cite{DillGPQXZ22}: without a heap, the verification conditions stay first-order and the solver needs no frame reasoning.
Running \WP on the eliminated form is a second use of the same transformation that, to our knowledge, has not been described before, and it avoids the following problems.
A predicate transformer over programs with references has to carry a model of locations: |*r| means the contents of whatever |r| points to, an assignment through one reference may change the value read through another, and the transformer needs either an explicit heap with separation or ownership reasoning to keep the substitutions sound~\cite{RustHorn,Aeneas}.
Its output is then stated over that model --- heap fragments, prophecy variables, ownership predicates --- and has to be translated before a user can read it.
On the eliminated form none of this arises.
Every reference is a path to a named root, so substitution is exact by construction; the only aliasing that survives is between global resources whose addresses may coincide, and it surfaces as an explicit condition such as |from == to| above.
The derived predicates mention only parameters, results, fields, and resources, which is the vocabulary of the specification language, so the reverse mapping of Sec.~\ref{sec:reverse} reduces to renaming.
And because inference and verification run on the same pipeline, every inferred contract is checked under the semantics that produced it; the instrumented contracts of callees and the loop invariants are available in-line to both.
The whole analysis is about nine thousand lines of Rust on top of the existing pipeline, and adds no new semantics to the prover.
Whether the same holds for a prophecy-based elimination, which is designed for \MVP but not implemented, is open; Sec.~\ref{sec:discussion} returns to it.

%%% Local Variables:
%%% mode: latex
%%% TeX-master: "main"
%%% End:

% !TEX root = main.tex
\Section{The Agent Side}
\label{sec:skills}

The agentic side of the workflow is delivered as a plugin to a general-purpose coding agent: skills that instruct it, hooks that check generated Move as it is written, and the Move toolchain exposed as tools it can call. The setup is at~\Anon{\cite{MoveFlow}}{\cite{MoveFlowAnon}}; we summarize the design.

\Paragraph{Inference tactics}
Inference runs in one of three \emph{tactics}, which differ in whether the \WP tool is available and who decides when to use it.
\begin{Itemize}
\item \emph{agent-only:} \WP is unavailable. The contract and any loop invariants are derived directly from the implementation and the contracts of its dependencies, checked as the work proceeds.
\item \emph{hybrid-guided:} the workflow of Sec.~\ref{sec:example}, prescribed as a fixed order --- run \WP over the scope, take one warned function and write an invariant from the loop-invariant evidence (Sec.~\ref{sec:diagnostics}), rerun \WP on that function alone until the warning is gone, simplify what \WP derived, then check.
\item \emph{hybrid-flexible:} \WP is available as one inference pass among others, and the agent decides whether and when to call it and how to combine it with its own reasoning.
\end{Itemize}

\Paragraph{Skills}
A tactic selects the \emph{plan} --- the sequence of tasks and which of them consult \WP --- and nothing else. All other guidance is shared reference material, identical in all three tactics: the specification language, the rules for writing and editing specs, the verification sub-workflow, and the toolchain's capabilities and limits. Only the description of the \WP tool itself is withheld where the tool is absent; the underlying idea of backward reasoning, and what a loop without an invariant costs, is background every tactic receives.

Three areas of that guidance carry most of the weight.
\begin{Itemize}
\item \emph{Loop invariants:} when to introduce recursive helper functions for iterative quantities, and when to prefer state-level invariants on resources --- assumed by the prover at every call site, including each loop iteration --- over per-loop inductive arguments.
\item \emph{Simplification:} discard vacuous conditions; replace unbounded quantifiers, a frequent source of SMT timeouts, with bounded or closed-form equivalents; and require trigger annotations on the surviving ones.
\item \emph{Timeout resolution:} a hierarchy of escalations --- strengthen state-level invariants, then introduce spec helpers, then lemmas, then explicit proof scripts --- with a strict prohibition on weakening abort or postconditions to silence a timeout, since this gives up a real property only to make the proof succeed.
\end{Itemize}

%%% Local Variables:
%%% mode: latex
%%% TeX-master: "main"
%%% End:

% !TEX root = main.tex
\Section{Evaluation}
\label{sec:evaluation}

The evaluation asks what the mechanical component adds to an agent, in completeness and in cost, and whether it changes which model is needed.
\emph{Arm against arm} (Sec.~\ref{sec:eval-arms}) holds the model fixed and varies the tactic: \textit{agent-only}, \textit{hybrid-guided}, and \textit{hybrid-flexible}, run under the same conditions.
\emph{Model against model} (Sec.~\ref{sec:eval-models}) repeats the design under a frontier model and a mid-tier model of the same family and asks whether, with \WP, the cheaper model reaches the results of the frontier model.
Both rounds share the corpus, the protocol, the scoring, and the cost accounting described first.

\SubSection{Setup}
\label{sec:eval-setup}

\Paragraph{Corpus}
The corpus is drawn from \Anon{Etna, the private Move code of the perpetual trading system Decibel}{the private Move code of a production perpetual trading system}: 23 targets from it, two from a public Move library that carry no upstream specification, and two written for the corpus.
A private codebase matters here: the model cannot have seen these functions, or any specification of them, during training, so a contract has to be inferred rather than recalled.
Every candidate is classified by the features (strata) its contract has to account for --- loops and accumulators, early returns and breaks, arithmetic overflow and underflow, casts and truncation, mutable references and in-place permutation, function values, quantifiers, composition through a callee's contract, and others, 38 in all.
A target is admitted only if it proves within seconds against a hand-written reference specification, and before the round three \emph{mutants} are authored per target --- patches to the implementation, one per obligation the contract has to pin --- each shown to be killed by the reference.
Twenty targets are selected to cover all 38 features across 15 modules, 10 of them with loops; three are deliberately guessable as controls, and five near-duplicates are held back.
With 3 arms and 4 replicates, a model round has \textbf{240 cells}.

\Paragraph{Arms and models}
The three tactics of Sec.~\ref{sec:skills} run as arms that differ only in the rendered plugin: the reference material is byte-identical, and in agent-only the \WP tool is absent rather than discouraged.
Two rounds use two models of one family through the same runtime, the Codex CLI at high effort: Sol~5.6, the frontier model, and Terra~5.6, the mid-tier model priced at about half of Sol's list.
Each session has a budget of 100--180k output tokens, one hour of wall time, and 40 seconds of solver time per verification condition.

\Paragraph{Scoring}
A cell reaches \emph{operational success} when a judge outside the agent's workspace confirms that the package compiles, the requested scope verifies, no weakening construct is present --- |pragma verify = false|, partial abort contracts, vacuous postconditions, invented preconditions --- the bytecode is unchanged, and the required contract categories are present.
The agent's own claim is never a score, and the agent receives only this acceptance feedback: the mutants are withheld until generation has ended, and a contract is never handed back for strengthening.
Operational success is necessary but weak: a vaguer contract passes it more easily.
Completeness --- postconditions as strong as the behavior --- has no deductive check in general; \WP guarantees it only given complete loop invariants and callee contracts, and warns when it cannot, a mechanism the agent-only arm lacks.
The second level is therefore \emph{mutation adequacy}: after the round, each withheld mutant is applied to the implementation and the prover rerun against the inferred specification.
A cell is \emph{complete} when it killed every mutant and \emph{failed} when one survived it; a cell whose scoring stayed inconclusive after one infrastructure retry is \emph{unmeasured} and counts in neither.
Every one of the 480 cells reached operational success, so the outcome of a cell is decided by the withheld mutants alone.

\Paragraph{Cost}
Costs are in US dollars at each model's published list prices, recomputed from the recorded token counters of every model request rather than taken from the harness's estimate; startup requests that do no task work are excluded.
Per million tokens, Sol~5.6 is \$4 uncached input, \$0.40 cached input, and \$20 output; Terra~5.6 is \$2, \$0.20, and \$12.
Thinking counts as output for both.
Means are task-weighted, and the 95\% intervals are percentile bootstraps that hold the 20 tasks fixed and resample the four replicate blocks within each task with the three arms kept together (100,000 draws); they describe this corpus, not a population of tasks.

\SubSection{Arm against Arm}
\label{sec:eval-arms}

Table~\ref{tab:summary} summarizes both rounds, Fig.~\ref{fig:cells} shows every cell, and Tables~\ref{tab:sol} and~\ref{tab:terra} in the appendix give every target.
Under both models the hybrids complete more cells than agent-only: 77 and 77 against 74 under Sol~5.6, and 79 and 78 against 74 under Terra~5.6.
All 16 failures occur on two targets.
|QP-part-025|, the in-place partition, is the one target whose contract has to characterize a whole data structure --- that the result is a rearrangement of the input --- and it defeats every arm under both models, hybrid-guided least under Terra~5.6 (1 cell) and most under Sol~5.6 (3 cells).
|OV-order-006| is the control that never aborts, because its |&&| short-circuits past the guards that would; agent-only stated an abort condition anyway in all four Sol~5.6 cells and two Terra~5.6 cells, and no hybrid cell missed it, since \WP derives the absence of an abort exactly.

\begin{table}[t]
\centering
\caption{Both rounds, per arm: mean cost in US dollars per cell, the hybrid/agent-only cost ratio with its 95\% bootstrap interval, cells that killed every withheld mutant, cells a mutant survived, cells whose scoring stayed inconclusive, mean input/output tokens in thousands, and mean wall time per cell in seconds.}
\label{tab:summary}
\scriptsize
\setlength{\tabcolsep}{2.5pt}
% Generated by eval_cells.py; do not edit.
\begin{tabular}{llrlrrrrr}
\toprule
model & arm & cost & vs. agent [95\% CI] & complete & failed & unmeas. & I/O k & wall s \\
\midrule
Sol~5.6 & agent-only & \$0.333 & --- & 74/80 & 4 & 2 & 250/4.9 & 119 \\
 & hybrid-guided & \$0.343 & 1.03 [0.94, 1.14] & 77/80 & 3 & 0 & 263/5.3 & 125 \\
 & hybrid-flexible & \$0.362 & 1.09 [1.00, 1.18] & 77/80 & 1 & 2 & 292/5.1 & 120 \\
\midrule
Terra~5.6 & agent-only & \$0.235 & --- & 74/80 & 5 & 1 & 372/6.4 & 124 \\
 & hybrid-guided & \$0.188 & 0.80 [0.73, 0.87] & 79/80 & 1 & 0 & 281/5.0 & 102 \\
 & hybrid-flexible & \$0.207 & 0.88 [0.78, 0.99] & 78/80 & 2 & 0 & 342/5.2 & 103 \\
\bottomrule
\end{tabular}

\end{table}

\begin{figure}[t]
\centering
% scatter/classes must sit on the plot, not in an axis style: applied through a
% style it accumulates per point and exhausts TeX's memory.
\newcommand{\cellplot}[1]{%
  \addplot[only marks, scatter, scatter src=explicit symbolic, mark size=1.25pt,
    scatter/classes={
      ok0={mark=*, ModelSol}, fail0={mark=o, ModelSol, thick}, unm0={mark=x, CellMuted, thick, mark size=1.8pt},
      ok1={mark=*, ModelTerra}, fail1={mark=o, ModelTerra, thick}, unm1={mark=x, CellMuted, thick, mark size=1.8pt}}]
    table[x index=1, y expr={\thisrowno{0} + (\thisrowno{2} - 0.5) * 0.36}, meta index=3] {#1};}
\begin{tikzpicture}
\begin{groupplot}[
  group style={group size=3 by 1, horizontal sep=2.5mm, yticklabels at=edge left},
  scale only axis, width=2.85cm, height=8.4cm,
  xmode=log, xmin=0.05, xmax=1.6, log ticks with fixed point,
  xtick={0.1,0.3,1}, xticklabels={0.1,0.3,1},
  minor xtick={0.05,0.2,0.5}, minor tick style={draw=none},
  ymin=-0.6, ymax=19.6, ytick={0,...,19}, y dir=reverse,
  ymajorgrids, grid style={CellGrid, thin},
  axis line style={CellAxis}, tick style={CellAxis},
  tick label style={font=\scriptsize},
  title style={font=\small, yshift=-1ex}]
\nextgroupplot[title=agent-only, yticklabels/.expanded={\CellTaskLabels}, yticklabel style={font=\ttfamily\scriptsize}]
\cellplot{tables/agent_only-cells.dat}
\nextgroupplot[title=hybrid-guided, xlabel={cost per cell (US\$)}, xlabel style={font=\scriptsize},
  legend style={at={(0.5,-0.11)}, anchor=north, font=\scriptsize, legend columns=-1, draw=none, /tikz/every even column/.append style={column sep=0.6em}},
  legend image post style={mark size=1.4pt}]
\cellplot{tables/hybrid_guided-cells.dat}
\addlegendimage{only marks, mark=*, ModelSol} \addlegendentry{Sol~5.6}
\addlegendimage{only marks, mark=*, ModelTerra} \addlegendentry{Terra~5.6}
\addlegendimage{only marks, mark=o, black, thick} \addlegendentry{mutant survived}
\addlegendimage{only marks, mark=x, CellMuted, thick, mark size=2pt} \addlegendentry{unmeasured}
\nextgroupplot[title=hybrid-flexible]
\cellplot{tables/hybrid_flexible-cells.dat}
\end{groupplot}
\end{tikzpicture}
\caption{Every cell of both rounds: cost in US dollars on a log scale, one panel per tactic, one row per target, and within each row the four Sol~5.6 replicates above the four Terra~5.6 replicates. A filled dot is a complete cell, a hollow one a cell a withheld mutant survived, a cross a cell whose scoring stayed inconclusive.}
\label{fig:cells}
\end{figure}

The effect on cost differs between the models.
Under Sol~5.6, hybrid-guided costs 1.03$\times$ agent-only [0.94, 1.14] and hybrid-flexible 1.09$\times$ [1.00, 1.18] in the task-weighted mean; both intervals reach parity, and the hybrids are cheaper on only 8 and 9 of the 20 targets.
The hybrids are cheaper on the small arithmetic contracts --- |VS-contrib-003|, |VS-shares-002|, |PM-curve-027| at about 0.6$\times$ --- where \WP's contract is already complete and the agent only has to simplify it, and more expensive on the targets that need invention beyond the loop invariant: |SM-select-022| at 1.75$\times$ under hybrid-guided, whose contract needs a recursive specification function over a function value, and the quantified-abort target |TL-lev-020| at 1.3--1.6$\times$.
On those targets the derived \WP output is large and is re-read on every turn.
The hybrids process a twentieth to a sixth more input per cell (263k and 292k against 250k) at similar output, and this additional input, billed at cache-read prices, accounts for the difference.

Under Terra~5.6 the hybrids are cheaper: hybrid-guided costs 0.80$\times$ agent-only [0.73, 0.87] and hybrid-flexible 0.88$\times$ [0.78, 0.99], cheaper on 14 of the 20 targets each, and both hybrids also finish faster (102 and 103 seconds against 124).
The direction of the token counts reverses: the hybrids read \emph{less} than agent-only (281k and 342k against 372k) and write less (5.0k and 5.2k against 6.4k), because the unaided model searches --- more verification calls, more repair rounds --- where the hybrids start from an exact contract.
The most expensive hybrid targets are the same three as under Sol~5.6, |SM-select-022|, |QP-part-025|, and |TL-lev-020|, at 1.2--1.3$\times$.

\SubSection{Model against Model}
\label{sec:eval-models}

Fig.~\ref{fig:cells} puts the two models side by side within each row, and the comparison is uniform: Terra~5.6 is cheaper than Sol~5.6 on 19 of the 20 targets in every arm, by 0.71$\times$ under agent-only and 0.55$\times$ and 0.57$\times$ under the hybrids in the task-weighted mean, at 231 against 228 complete cells over the round.
The exception is |BA-base-012|, the lemma target, where Terra~5.6 costs 1.2--1.5$\times$ Sol~5.6 because it needs more attempts to find the monotonicity lemma.
Since Terra's price list is about half of Sol's, the agent-only ratio of 0.71 implies that Terra spends about 40\% \emph{more} tokens per cell to reach the same result, while the hybrid ratios near 0.55 imply that with \WP it spends about the same: the tool closes the token gap between the two models.
In the figure, the replicate spread is about the same for both models in every arm, a factor of about 1.6 between the cheapest and the most expensive replicate of a typical target.
What changes with the tactic is how far the Terra replicates lie to the left of the Sol replicates: little under agent-only, and clearly on nearly every row under hybrid-guided.

Together, the two rounds show two effects of \WP.
The first is on completeness and does not depend on the model: under both, the hybrids complete more cells, never miss the control whose abort condition \WP derives exactly, and fail only on the one target whose contract needs a permutation argument that \WP cannot supply either.
The second is on cost and depends on the model.
\WP supplies exact contracts, but its output also enlarges the context the model re-reads on every turn.
Sol~5.6 finds the invariants unaided, so it pays for the additional context without saving elsewhere (1.03--1.09$\times$, parity within the intervals).
Terra~5.6 would otherwise search, and the savings outweigh the context (0.80$\times$); with \WP, it completes as many cells as Sol~5.6 at 55\% of its cost.
With four replicates and failures concentrated on two targets, the completeness differences are descriptive; the Terra~5.6 cost intervals exclude parity, the Sol~5.6 ones include it.

%%% Local Variables:
%%% mode: latex
%%% TeX-master: "main"
%%% End:

% !TEX root = main.tex
\Section{Discussion and Conclusion}
\label{sec:discussion}

\SubSection{Interpreting the Results}

The evaluation has three results. An agent without mechanical help completes 99\% of its
cells under the frontier model and 93\% under the mid-tier model. \WP does not change
completeness, but it lowers the cost wherever it derives the contract by itself. And the
mid-tier model nearly reaches the results of the frontier model, at a cost that follows
its lower price.

\Paragraph{Where \WP pays off}
On targets without loops, \WP needs no help: it derives the complete contract, and the
agent only checks it, at 36--60\% of the cost of writing the contract directly. Half of
this corpus is of this kind, including four of the five compositions over callee
contracts and eight of the nine targets with global state. On loop targets the agent still has to find the
invariant, and \WP's quantified output adds context that is re-read on every turn, so
the hybrids cost slightly more there. With \WP, the first candidate check is accepted more
often, most for the mid-tier model (from 37\% to 63--65\% of cells).

\Paragraph{Where it does not}
The partition needs a permutation argument that \WP cannot supply, and every tactic fails
on it, the mid-tier model more often. The two scans that leave their loop early fail only
in the hybrid tactics: the agent accepted the contract that \WP derived from its prefix
invariants once it verified, and a mutant on the exit path survived. A derived contract
that verifies can thus still be incomplete, and the tool's diagnostics, which report a
missing invariant, do not detect this case.

\Paragraph{Mid-tier and frontier models}
Both models spend about the same number of tokens per cell, so the cost of the mid-tier
model follows its price. Apart from the permutation target, it handles the loop
invariants, compositions, and global state of the corpus as well as the frontier model.
We expect the cost advantage of \WP to extend to cheaper models: for a local model, the
context \WP adds costs latency rather than money, and a model trained to read \WP output
should need fewer rounds with the tool.

\Paragraph{Carrying over to prophecies}
The reverse mapping of Sec.~\ref{sec:reverse} relies on the elimination being
\emph{path-based}: an eliminated reference is still a path to a location the source can
name, so the derived predicate has a source-level spelling. A \emph{prophecy}-based
elimination in the style of RustHorn and Creusot~\cite{RustHorn}, in which a mutable
reference is a pair of its current value and a prophecy variable for its final value,
pinned when the borrow ends, has been designed for \MVP but not implemented; it would
need no borrow paths and would handle borrows through function values, which the
path-based model cannot. A contract is model-agnostic, so inference could keep running
on the path-based form regardless. Whether \WP could run on the prophecy form itself
and still be carried back is open. For a |&mut| parameter the prophecy \emph{is} the
post-state value and has the spelling |p|; for a reborrow through a function value the
final value of the borrow has no source-level counterpart at all, and a predicate over it
would have to be existentially closed or restated through the callee's contract. We do
not know yet whether that always yields a readable contract.

\Paragraph{Open questions}
The evaluation leaves the questions of scale open. Its targets are functions, not
modules; its corpus comes from one codebase; its two models are of one family; and its
cost measure is dollars under their price schedules. We plan to apply inference to the
parts of the Aptos framework~\cite{AptosFrameworkBook} that already carry hand-written
specifications and to compare the inferred contracts with them --- a comparison this
corpus cannot make, since its targets were selected for having no specification.

\SubSection{Related Work}

Specification inference predates large language models --- Daikon~\cite{DaikonTSE01} learned likely invariants from program traces --- but the techniques that followed stayed tied to particular property classes and source languages, and the advent of LLMs has revived interest in the area.
Pei et al.~\cite{PeiEtAl23} showed that fine-tuned LLMs can predict program invariants statically at quality competitive with dynamic analysis.
Kamath et al.~\cite{KamathEtAl23} pair an LLM that \emph{proposes} loop invariants with a symbolic checker that \emph{validates} them, retrying on failure --- a propose-and-validate loop that recurs in much of the subsequent work.
For Verus, AlphaVerus~\cite{AlphaVerus} bootstraps verified code by self-improving translation with tree search over verifier feedback, and VeriStruct~\cite{VeriStruct26} synthesizes Verus-style specifications and proofs for Rust functions, modules, and data structures through planner-orchestrated phases.
For smart contracts in Solidity, PropertyGPT~\cite{PropertyGPT25} retrieves human-written Certora properties as input examples for specification generation.
Most directly related to our work, MSG~\cite{MSG25} targets the same Move specification language through a skill-based agentic design that consumes verifier feedback beyond a binary success/failure verdict.
Across these systems the LLM remains responsible for the entire specification; the verifier serves as a binary oracle or, in MSG and our setup, as a feedback channel, but never as a contributor of mechanically derived content.
A separate line of work, exemplified by DafnyPro~\cite{DafnyPro26} and AutoVerus~\cite{AutoVerus25}, tackles the complementary problem of synthesizing the auxiliary proof annotations a verifier needs when the specification is already given. Our plugin can also synthesize proof hints (Sec.~\ref{sec:pow}), but this is not the focus of this paper.

Reference elimination, the transformation underlying our \WP analysis, was introduced in~\cite{DillGPQXZ22}, and equivalents exist for Rust.
Aeneas~\cite{Aeneas} compiles Rust to a pure functional model in a target prover (Lean, F*, Coq, HOL4) via forward and backward functions, while RustHorn~\cite{RustHorn} encodes a mutable borrow as a pair of its current value and a prophecy for its value at the end of the borrow, reducing verification to constrained Horn clauses.
Neither infers specifications --- in Aeneas the user writes them in the target prover, against the translated model --- and both consume the reference-free form only in the forward direction.
Our use of the transformation is the opposite: \WP runs on the reference-free form, and the conditions it derives are carried \emph{back} through the elimination.
To our knowledge this inversion has not been described before.
Boogie~\cite{boogie,boogieIVL}, the intermediate verification language that \MVP translates to, also computes weakest preconditions of the procedures it is given~\cite{BL05} and passes them to the SMT solver; the difference is that Boogie's \WP is not intended for external consumption, whereas ours has to be read, which is reflected in the reverse mapping and in the effort spent on simplification.

Two aspects thus distinguish our setup from the work above.
First, to our knowledge, no prior LLM-based spec-inference system uses a sound mechanical baseline alongside the LLM.
\WP itself produces a substantial part of the inferred specification --- abort conditions, modifies clauses, and consequence-style postconditions --- so the AI is left only with what \WP cannot recover, and the two streams are validated together by the Move Prover.
Second, the workflow is delivered as a plugin to a general-purpose, interactive coding agent via skills, edit-time hooks, and an MCP server, rather than as a custom-built tool: the agent participates in the loop and the user can intervene at any point.

\SubSection{Conclusion}
\label{sec:conclusion}

We have presented a specification inference approach for the Move Prover that runs \WP on reference-eliminated bytecode, where backward propagation is substitution and the derived contracts can be mapped back to the references of the source, and we have combined it with a coding agent that supplies what \WP cannot: loop invariants, simplification, and proof hints. The prover checks both.
Reference elimination was introduced for verification; this work shows that it serves inference as well, and the approach should carry over to verifiers for Rust-like languages that eliminate references along paths.
A controlled comparison on 26 targets, 24 of them from public code of the Aptos framework without specifications --- three tactics, four replicates each, two models --- found that a pure agent produces a complete specification in 93--99\% of its cells; that \WP leaves completeness unchanged but lowers the cost by 11--21\%, all of it on targets without loops, where the hybrids cost 36--60\% of the agent alone; and that the mid-tier model costs 53--59\% of the frontier model and falls behind it mainly on one target.
We expect local and specialized models, for which the additional context costs least, to benefit further from \WP.

%%% Local Variables:
%%% mode: latex
%%% TeX-master: "main"
%%% End:


\subsubsection*{Data Availability Statement.}
The round archives behind every table and figure --- per-cell costs, token counters, wall times, and mutant-gate outcomes --- and the scripts that produce them from the archives are published with the inference plugin at~\Anon{\cite{MoveFlow}}{\cite{MoveFlowAnon}}.
The corpus targets are extracted from a private codebase and are not published; the two targets written for the corpus are included.

\bibliography{biblio}

\appendix
\Section{Per-Target Results}
\label{sec:appendix}

Tables~\ref{tab:sol} and~\ref{tab:terra} give, for each target and arm, the mean cost per cell over the four replicates, mean input/output tokens in thousands (I/O k), the hybrid mean cost as a percentage of agent-only (vs. agent), and whether the specification killed every withheld mutant in all four replicates (\Checkmark) or a mutant survived (\XSolidBrush, subscripted with the count where more than one replicate failed); $\circ$ marks a cell whose scoring stayed inconclusive.
Input counts include uncached, cached, and cache-creation tokens.

\begin{sidewaystable}
\centering
\caption{Sol~5.6 by target. Summed over the 20 targets the arms cost \$26.65, \$27.44, and \$28.97.}
\label{tab:sol}
\footnotesize
\setlength{\tabcolsep}{2.2pt}
% Generated by eval_cells.py from corpus3.2-run11-codex-sol56-high-exact-cost.tar.gz; do not edit.
\begin{tabular}{lrrcrrrcrrrc}
\toprule
& \multicolumn{3}{c}{agent-only} & \multicolumn{4}{c}{hybrid-guided} & \multicolumn{4}{c}{hybrid-flexible} \\
\cmidrule(lr){2-4} \cmidrule(lr){5-8} \cmidrule(lr){9-12}
target & cost & I/O k & ok & cost & vs. agent & I/O k & ok & cost & vs. agent & I/O k & ok \\
\midrule
|BA-base-012| & \$0.535 & 449/9.7 & \Checkmark & \$0.554 & 104\% & 398/9.7 & \Checkmark & \$0.757 & 141\% & 677/14.3 & \Checkmark \\
|BK-bucket-016| & \$0.331 & 329/4.1 & \Checkmark & \$0.408 & 123\% & 403/4.8 & \Checkmark & \$0.305 & 92\% & 257/3.3 & \Checkmark \\
|LP-price-021| & \$0.280 & 199/3.5 & \Checkmark & \$0.257 & 92\% & 141/3.7 & \Checkmark & \$0.272 & 97\% & 185/4.8 & \Checkmark \\
|MD-median-015| & \$0.152 & 100/1.5 & \Checkmark & \$0.191 & 125\% & 107/1.6 & \Checkmark & \$0.195 & 128\% & 137/1.5 & \Checkmark \\
|MM-min-013| & \$0.307 & 208/2.5 & \Checkmark & \$0.288 & 94\% & 219/3.9 & \Checkmark & \$0.279 & 91\% & 189/3.6 & \Checkmark \\
|OV-order-006| & \$0.187 & 120/2.3 & \XSolidBrush$_{4}$ & \$0.190 & 102\% & 104/1.3 & \Checkmark & \$0.233 & 124\% & 125/1.6 & \Checkmark \\
|PM-curve-027| & \$0.229 & 137/2.8 & \Checkmark & \$0.147 & 64\% & 91/1.2 & \Checkmark & \$0.194 & 84\% & 141/1.4 & \Checkmark \\
|QP-part-025| & \$0.618 & 473/12.7 & $\circ$ & \$0.708 & 115\% & 677/13.6 & \XSolidBrush$_{3}$ & \$0.708 & 115\% & 610/12.9 & \XSolidBrush \\
|SM-select-022| & \$0.450 & 304/9.2 & \Checkmark & \$0.787 & 175\% & 480/19.0 & \Checkmark & \$0.473 & 105\% & 320/9.9 & $\circ_{2}$ \\
|TL-lev-020| & \$0.271 & 208/3.7 & \Checkmark & \$0.349 & 129\% & 287/5.0 & \Checkmark & \$0.432 & 159\% & 397/5.5 & \Checkmark \\
|TR-cancel-026| & \$0.347 & 257/5.0 & $\circ$ & \$0.354 & 102\% & 294/5.7 & \Checkmark & \$0.360 & 104\% & 266/5.3 & \Checkmark \\
|TR-discard-011| & \$0.299 & 213/3.5 & \Checkmark & \$0.329 & 110\% & 257/4.5 & \Checkmark & \$0.360 & 120\% & 271/4.1 & \Checkmark \\
|TR-order-010| & \$0.226 & 172/2.9 & \Checkmark & \$0.299 & 132\% & 262/4.2 & \Checkmark & \$0.318 & 141\% & 297/4.3 & \Checkmark \\
|TS-trial-019| & \$0.362 & 220/5.3 & \Checkmark & \$0.236 & 65\% & 134/2.8 & \Checkmark & \$0.234 & 64\% & 156/2.3 & \Checkmark \\
|UC-credits-008| & \$0.391 & 267/4.8 & \Checkmark & \$0.324 & 83\% & 251/4.9 & \Checkmark & \$0.333 & 85\% & 287/4.5 & \Checkmark \\
|VS-contrib-003| & \$0.307 & 280/3.6 & \Checkmark & \$0.177 & 58\% & 122/1.3 & \Checkmark & \$0.247 & 80\% & 155/1.7 & \Checkmark \\
|VS-fees-001| & \$0.348 & 276/5.1 & \Checkmark & \$0.371 & 107\% & 305/6.7 & \Checkmark & \$0.504 & 145\% & 422/7.2 & \Checkmark \\
|VS-redeem-004| & \$0.535 & 446/8.9 & \Checkmark & \$0.494 & 92\% & 497/7.5 & \Checkmark & \$0.614 & 115\% & 635/9.5 & \Checkmark \\
|VS-shares-002| & \$0.274 & 200/3.0 & \Checkmark & \$0.171 & 62\% & 99/1.6 & \Checkmark & \$0.222 & 81\% & 171/1.9 & \Checkmark \\
|WU-consume-023| & \$0.211 & 136/2.8 & \Checkmark & \$0.228 & 108\% & 137/3.0 & \Checkmark & \$0.202 & 96\% & 151/2.4 & \Checkmark \\
\midrule
\textbf{all 20} & \$0.333 & 250/4.9 & 74/80 & \$0.343 & 103\% & 263/5.3 & 77/80 & \$0.362 & 109\% & 292/5.1 & 77/80 \\
\bottomrule
\end{tabular}

\end{sidewaystable}

\begin{sidewaystable}
\centering
\caption{Terra~5.6 by target. Summed over the 20 targets the arms cost \$18.84, \$15.01, and \$16.60.}
\label{tab:terra}
\footnotesize
\setlength{\tabcolsep}{2.2pt}
% Generated by eval_cells.py from corpus4-run8-codex-terra56-high.tar.gz; do not edit.
\begin{tabular}{lrrcrrrcrrrc}
\toprule
& \multicolumn{3}{c}{agent-only} & \multicolumn{4}{c}{hybrid-guided} & \multicolumn{4}{c}{hybrid-flexible} \\
\cmidrule(lr){2-4} \cmidrule(lr){5-8} \cmidrule(lr){9-12}
target & cost & I/O k & ok & cost & vs. agent & I/O k & ok & cost & vs. agent & I/O k & ok \\
\midrule
|keep_alive| & \$0.273 & 585/4.2 & \Checkmark & \$0.060 & 22\% & 69/0.4 & \Checkmark & \$0.102 & 37\% & 176/0.9 & \Checkmark \\
|is_order_valid| & \$0.269 & 574/3.7 & \Checkmark & \$0.062 & 23\% & 67/0.4 & \Checkmark & \$0.110 & 41\% & 154/0.9 & \Checkmark \\
|available_transaction_queue_capacity| & \$0.113 & 182/1.9 & \Checkmark & \$0.059 & 52\% & 71/0.4 & \Checkmark & \$0.085 & 75\% & 124/0.7 & \Checkmark \\
|can_execute_with_timelock| & \$0.239 & 470/4.4 & \XSolidBrush & \$0.059 & 25\% & 66/0.4 & \Checkmark & \$0.098 & 41\% & 165/0.9 & \Checkmark \\
|is_taker_order| & \$0.136 & 227/1.9 & \Checkmark & \$0.061 & 45\% & 66/0.4 & \Checkmark & \$0.088 & 64\% & 117/0.7 & \Checkmark \\
|refill| & \$0.168 & 267/3.9 & \Checkmark & \$0.060 & 35\% & 67/0.4 & \Checkmark & \$0.096 & 57\% & 141/0.9 & \Checkmark \\
|request| & \$0.100 & 151/1.6 & \Checkmark & \$0.059 & 59\% & 73/0.5 & \Checkmark & \$0.122 & 122\% & 171/1.9 & \Checkmark \\
|validate_enough_stake| & \$0.217 & 358/5.1 & \Checkmark & \$0.141 & 65\% & 239/2.4 & \Checkmark & \$0.121 & 56\% & 223/1.9 & \Checkmark \\
|find_min_stake_required| & \$0.175 & 321/2.3 & \Checkmark & \$0.083 & 47\% & 114/1.1 & \Checkmark & \$0.118 & 68\% & 201/1.7 & \Checkmark \\
|update_config| & \$0.159 & 289/2.7 & \Checkmark & \$0.065 & 41\% & 69/0.4 & \Checkmark & \$0.105 & 66\% & 175/0.9 & \Checkmark \\
|match_order_and_get_next_from_bulk_order| & \$0.257 & 517/4.8 & \Checkmark & \$0.065 & 25\% & 68/0.4 & \Checkmark & \$0.086 & 34\% & 127/0.8 & \Checkmark \\
|reinsert_order_into_bulk_order| & \$0.192 & 340/3.2 & \Checkmark & \$0.057 & 30\% & 66/0.4 & \Checkmark & \$0.079 & 41\% & 118/0.7 & \Checkmark \\
|new_bulk_order_request_with_sanitization| & \$0.173 & 263/3.8 & \Checkmark & \$0.059 & 34\% & 75/0.5 & \Checkmark & \$0.095 & 55\% & 141/0.8 & \Checkmark \\
\midrule
|validate_scheme| & \$0.142 & 224/3.0 & \Checkmark & \$0.173 & 122\% & 292/3.5 & \Checkmark & \$0.202 & 142\% & 361/4.3 & \Checkmark \\
|upsert_provider_jwks| & \$0.310 & 543/8.1 & $\circ$ & \$0.220 & 71\% & 346/5.1 & \XSolidBrush$_{2}$ & \$0.335 & 108\% & 603/9.1 & \XSolidBrush \\
|get_pending_transactions| & \$0.306 & 663/5.1 & \Checkmark & \$0.325 & 106\% & 644/7.8 & \Checkmark & \$0.392 & 128\% & 850/9.2 & \Checkmark \\
|partition| & \$0.198 & 372/3.7 & \XSolidBrush$_{4}$ & \$0.476 & 240\% & 1129/11.3 & \XSolidBrush$_{2}$ & \$0.255 & 129\% & 489/5.8 & \XSolidBrush$_{4}$ \\
|select| & \$0.338 & 612/9.0 & $\circ$ & \$0.558 & 165\% & 1224/15.0 & $\circ$ & \$0.435 & 129\% & 906/10.3 & \Checkmark \\
|split_signature_bytes| & \$0.185 & 351/2.8 & \Checkmark & \$0.172 & 93\% & 339/2.6 & \Checkmark & \$0.193 & 105\% & 355/3.7 & \Checkmark \\
|new_tiers| & \$0.232 & 412/5.9 & \Checkmark & \$0.238 & 102\% & 394/5.5 & \Checkmark & \$0.320 & 138\% & 660/7.4 & \Checkmark \\
|validate_tiers| & \$0.195 & 381/3.4 & \Checkmark & \$0.260 & 133\% & 595/4.8 & \Checkmark & \$0.320 & 164\% & 745/5.7 & \Checkmark \\
|cancel_at_price_level| & \$0.192 & 314/4.3 & \Checkmark & \$0.194 & 101\% & 344/4.3 & \XSolidBrush$_{2}$ & \$0.188 & 98\% & 326/3.6 & \XSolidBrush \\
|discard_price_crossing_levels| & \$0.194 & 343/3.5 & \Checkmark & \$0.144 & 74\% & 220/2.5 & \Checkmark & \$0.184 & 95\% & 317/3.3 & \Checkmark \\
|validate_not_zero_sizes| & \$0.125 & 202/1.6 & \Checkmark & \$0.144 & 115\% & 225/1.7 & \Checkmark & \$0.140 & 112\% & 272/2.1 & \Checkmark \\
|validate_price_ordering| & \$0.159 & 255/2.2 & \Checkmark & \$0.144 & 91\% & 224/2.3 & \Checkmark & \$0.164 & 103\% & 293/2.8 & \Checkmark \\
|range_with_step| & \$0.216 & 331/5.4 & \Checkmark & \$0.242 & 112\% & 426/5.8 & \Checkmark & \$0.264 & 122\% & 489/6.7 & \Checkmark \\
\midrule
\textbf{all 26} & \$0.202 & 367/3.9 & 97/104 & \$0.161 & 79\% & 289/3.1 & 97/104 & \$0.181 & 89\% & 335/3.4 & 98/104 \\
\bottomrule
\end{tabular}

\end{sidewaystable}

\end{document}
